\documentclass{article}
\usepackage[utf8]{inputenc}
\usepackage[T1]{fontenc}
\usepackage{geometry}
\usepackage{graphicx}
\usepackage{float}
\usepackage{hyperref}
\usepackage[english,indonesian]{babel}

\title{Optimasi Pemrosesan Template LaTeX Menggunakan Pendekatan Berbasis AI}
\author{Academic Support Agent}
\date{\today}

\begin{document}
\maketitle

\begin{abstract}

\end{abstract}

\chapter{PENDAHULUAN}


\chapter{TINJAUAN PUSTAKA}


% Original content (commented out):
% \include{prelude}

\title{[Desarrollo de Software/Ingeniería de Software] Plantilla}

\begin{document}
	\include{coverpage}
	
    \section{Introducción}
	
    \section{Objetivos del proyecto}
    
    \subsection{Objetivos general}
    
    \subsection{Objetivos particulares}

    \section{Descripción de la problemática}
       
    \subsection{Motivación}
    
    \subsection{Definición del problema}
        
    \subsection{Estado del arte}
   
      Ejemplo de listado
      \begin{itemize}
      \setlength\itemsep{0em}
      \item \textbf{Algo} \cite{c1}
      \item \textbf{Otro algo} \cite{c2}
      \item \textbf{Un tercer algo} \cite{c3}
      \end{itemize}
       
    \section{Descripción de la solución propuesta}
    
    \subsection{Características de la solución}
	
    \subsection{Propósitos de la solución}

	\subsection{Alcances y limitaciones de la solución}
    
    \section{Metodología, herramientas y ambiente de desarrollo}
   
    \subsection{Metodología a usar}
  
    \subsection{Herramientas de desarrollo}
     Ejemplos de tablas:
     
     En la Tabla \ref{table:software}, se nombran el software y su respectiva función en el proyecto. 
	%%% Ver ejemplo en el apartado 5.3.2
    
    En la Tabla \ref{table:listaequipos}, se describe el hardware a utilizar en el proyecto. Cada dispositivo utilizado servirá para múltiples tareas, ya sea de desarrollo, documentación o pruebas del proyecto.
    %%% Ver ejemplo en el apartado 5.3.2
   
    \subsubsection{Software}
        
	    \begin{table}[!ht]
        \centering
        \caption{Software} \label{table:software}
        \begin{tabularx}{\textwidth}{|X|X|}
        	\hline
            \bf Uso de la Herramienta & {\bf Nombre} \\
            \hline
            {\bf Documentación} & {\LaTeX $2\varepsilon$} \newline 
                                  {Microsoft Office 2016} \newline
                                  {OpenOffice 4.0.1}\\
            \hline
            {\bf Sistema operativo} & {Windows 7 Ultimate (64 bits)} \newline
			 						  {Windows 10 Enterprise (64 bits)} \newline
             						  {Windows 10 Professional (64 bits)}\\
            \hline
            {\bf Lenguaje de Marcado} & {HTML 5} \newline
             							{CSS 3}\\ 
            \hline
            {\bf Lenguaje de Programación}& {Java SE Development Kit 7u80}\\
            
            \hline
            {\bf Base de Datos} & {Oracle 11g XE}\\
            \hline
            {\bf Herramienta de diseño Base de Datos} & {Oracle Designer 10.1.2.6 Release}\\
            \hline
            {\bf Editor de texto IDE} & {Notepad++ 6.7.9.2} \newline
                     					{Netbeans 8.1}\\
            \hline	
            \end{tabularx}       
        \end{table}  
    %\newpage    
    \subsubsection{Hardware}
    	
    \begin{table}[!ht]
    	\centering
        \caption{Hardware} \label{table:listaequipos}
        \begin{tabularx}{\textwidth}{|p{2cm}@{}|X|X|p{2cm}@{}|}
           \hline
           {\bf Equipo} & {\bf Sistema Operativo} & {\bf Procesador} & {\bf RAM}\\
           \hline
           {\bf PC} & {Windows 10 Enterprise (64 bits)} & {Intel® Core™ i5-4460 3.2GHz} & {8GB} \\
           \hline%nano
           {\bf Notebook 1} & {Windows 7 Ultimate (64 bits)} & {Intel® Core™ i3-380m 2.53GHz} & {6GB} \\
           \hline%nano
           {\bf Notebook 2} & {Windows 10 Home (64 bits)} & {Intel® Core™ i5-6200 2.3GHz} & {6GB} \\
           \hline%barbara
           {\bf Notebook 3} & {Windows 8 Professional } & {AMD® Turion™ X2 Dual-Core Mobile RM-72 2.10GHz} & {4GB} \\
           \hline
           {\bf Notebook 4} & {Windows 10 Enterprise (x64)} & {Intel® Core™ i7-6700HQ (2600 MHz - 3500 MHz)} & {16GB DDR4} \\
           \hline
        \end{tabularx}
    \end{table}
    
    \subsection{Ambiente de desarrollo}
	
    \subsubsection{Recursos materiales} %% Subsección acorde al recurso
    
	\subsubsection{Recursos humanos} %% Subsección acorde al recurso
    
    Los recursos humanos se pueden apreciar en la Tabla \ref{table:rrhh} con los cargos de cada integrante.
    %%%% MODELO DE TABLA
    %\vspace*{-\baselineskip} %%% ARREGLA EL SALTO DE LINEA EXTRA ENTRE TEXTO Y TABLA
    \begin{table}[!ht] %%% LA TABLA SE COLOCA EN DONDE SE ORIGINA, CON AJUSTES RESPECTIVOS
      \centering %%% CENTRADO
      \caption{Recursos humanos}\label{table:rrhh} %%% Leyenda y Etiquetado
      \begin{tabularx}{\textwidth}{|p{4cm}@{}|X|} %%% En este ejemplo, dos celdas con alineación a la izquierda
	      \hline 
    	  {\bf Nombre del integrante}     & {\bf Cargo en el proyecto} \\ \hline
      	  {Integrante 1} & {}                     \\ \hline
          {Integrante 2}      & {}                     \\ \hline
          {Integrante 3}      & {}                     \\ \hline
          {Integrante 4}  & {}                     \\ \hline
          \end{tabularx}
    \end{table}

    \section{Plan de trabajo}
    
    \subsection{Planificación de entrevistas}
	
    \subsection{Planificación del proyecto}
    
    \section{Conclusiones}
	
    \include{bibliography}
    
    \section{Apéndices}

\end{document}

\end{document}